\subsubsection{完整定价怎样把计算结果转为下界}
任务列主问题首先看到的是机型总工作量。两架飞机的工作量除以2给出平均界，但实际任务不能被分成任意小片段，在两条时间轴上的装填可能存在离散不平衡。因此装填割增加了任务长度分布的信息，覆盖约束保证物资数量，分支条件进一步约束整数聚合量，三者共同加强松弛。

另一方面，受限列池可能遗漏便宜任务，直接使用其最优值并不足以形成原问题下界。完整定价检查所有允许列的对偶不等式；Feillet等关于资源约束基本最短路的精确标签算法为此类状态扩展与支配设计提供了方法基础\cite{feillet2004}。本题的定价除了容量和能源，还要保留相关分支状态，使支配删除不会遗漏影响对偶可行性的列。

若全部允许列满足式\eqref{eq:q2_pricing_rc}，则对任一映射到该叶域的实际方案，有
\begin{equation}
\begin{aligned}
 \pi^{\mathsf T}d
 &\le\sum_g\mu_g\sum_{r:g(r)=g}\underline p_rx_r+\lambda^{\mathsf T}Bx\\
 &\le T\sum_gn_g\mu_g+\lambda^{\mathsf T}b
 \le T+\lambda^{\mathsf T}b.
\end{aligned}
\label{eq:q2_weak_dual_chain}
\end{equation}
移项即为$T\ge\pi^{\mathsf T}d-\lambda^{\mathsf T}b$。这个推导说明，完整定价保证的是证书中的对偶可行性，而不只是“列池不再变化”。各叶域分别得到有效界，再通过互补覆盖统一到原问题，才形成最终认证区间。

上界与下界可以沿不同路线持续改善：重构任务和时序降低可行上界，加入有效约束与加强定价提高下界。当前两者相距259.275062 s，量化了名义实例中尚可能存在的工期改进空间。可执行方案与独立界定相结合，使解质量不再仅由一次启发式输出的数值判断。
